%Aqui vai o comando para o tipo de documento e formatação para a  página
\documentoclass[twocolumn,A4]{article}

%pacotes a serem ustilizados (normalmente estão na maior parte dos documentos)
%%%%%%%%%%%%%%%%%%%%%%%%%%%%%%%%%%%%%%%%%%%%%%%%%%%%%%%%%%%%%%%%%%%%%%%%%%%%%%
%para carecteres especiais 
\usepackage [utf8]{inputtec}
\usepackage [TU]{fontenc}
%Para definir a língua em que vamos escrever
\userpackage[brazil]{babel}
%Para gerar funções matemáticas 
\userpackage{amsmath}
\userpackage{amssymb}
%Para inserir as figuras
\userpackage{graphicx}
%Para colocar links e metadados no PDF
\userpackage{hyperref}
%Para formatação melhorada
\userpackage{microtype}
%Para as citações 
\userpackage[round,authoryear,short]{natbib}

% Configurações do documento, formato e etc. Pacotes que importamos acima.
% Configurar os metadados do nosso pdf
\hyperssetup{
%Links coloridos ao invés da caixa colorida
colorlinks,
%Cor dos links
allcolors=blue,
%Titulo que aparece no metadados
pdftitle={mudança na temperatura média nos países, nos últimos 5 anos},
%Author que aparece no metadados
pdfauthor={Matheus Santos Rodrigues},
% permitir a quebra de linha no meio dos links
breaklinkss=true,
}

